\documentclass[../main.tex]{subfiles}

\begin{document}

\chapter*{Preface}

Hello, thank you for picking up this book. I have just finished writing it, except for this 
portion which I'm writing right now! It has been one epicly long journey of 9 weeks in my 
last quarter of college before I graduate. I wasn't even sure if I would be able to complete 
this by the end, but it seems just like my assignments, I finished this at the last minute.

To be completely honest, I am not an expert when it comes to computer science. After all, 
my degree is in environmental sciences (just kidding, I am a computer science major).
However, even though I don't know everything about every data structure, algorithm, 
design pattern, or whatever, I think that this book would be great for someone getting 
started in the world of SAT solving. These things that I am not an expert in are my passion,
and by doing this technical research, I hope to become a little less of ``not an expert'' in 
those fields.

We are not going to revolutionize SAT solving, in fact, we're barely going to scratch the 
surface in this book (it's a big world out there). What I aim to do in writing this is to 
help you understand how SAT solvers work and how they are implemented. I hope that this 
remains accessible and I'm not too confusing, however, I am just a \sout{girl} undergraduate.
Because we're covering how SAT solvers work and are implemented, if you are a student tasked 
with building one, don't use this to cheat. That would hurt my feelings. Instead use this to 
help your own understanding instead of copy-pasting.

Throughout this book we will build our own modern SAT solver from scratch, 
complete with all the bells and whistles. We will begin with the most basic 
algorithm imaginable, brute force, and gradually work our way through DPLL, CDCL, 
and beyond.

Throughout this journey we will encounter thought experiments, exercises, and a
handful of experiments intended to solidify our understanding. Any external resources will 
be cited (hopefully correctly), and in keeping with the goal of accessibility:
\begin{itemize}
    \item Questions that are explicitly asked will eventually be answered.
    \item Proofs will be explained in informal language before becoming formal.
    \item Code will be provided and linked through GitHub (or another hosting platform).
    \item Revisions to this book will be made whenever I discover mistakes or better explanations.
\end{itemize}

There is one more thing that deserves acknowledgement: some pieces of code were AI-assisted.
To be completely transparent, ``AI-assisted'' means that instead of spending days repeatedly 
banging my head against a wall trying to understand a bug or some subtle algorithmic detail, 
I occasionally asked an AI to explain concepts or help debug an implementation. Every word of this book, 
however, is my own. AI did not generate any of the explanatory text you will read here, 
although it certainly helped produce some of the code that accompanies it. I know that AI 
is very unpopular, and I do get that, however, even with literature and code I could find, it 
wasn't enough \textbf{for me} to figure out how everything is implemented from the ground up. 
That being said, I would still encourage you to read the other literature out there and look 
at code on your own, and perhaps you'll find me an idiot for needing AI assistance in the first 
place.

There are practical reasons for this. SAT solvers are difficult to implement correctly, 
and introducing subtle bugs is remarkably easy. In fact, before deciding to write this book, 
I had already attempted to build a SAT solver once before. I never managed to get a working CDCL 
implementation.

That failure, however, taught me a tremendous amount. So in a sense, we are learning
SAT solving together (isn't that comforting?). My previous mistakes have shown me many 
of the pitfalls, and my hope is that by documenting this journey I can help you avoid making the same ones.
So without further ado, let's begin. By the time we reach the end, I hope I will have 
taught you more than I myself learned.

\end{document}